% !TeX program = lualatex
% halo:begin
% {
%   "version": 1,
%   "publish": true,
%   "course": "mathematical-analysis-i",
%   "section": "2.2",
%   "title": "数学分析（一）2.2：收敛数列的性质",
%   "categories": ["study-notes"],
%   "tags": ["mathematics"],
%   "excerpt": "收敛数列极限的唯一性、有界性、保号性及最终区间性质。"
% }
% halo:end
% 来源：IMG_1369.HEIC、IMG_1370.HEIC。
\RequirePackage{iftex}
\RequireLuaTeX
\documentclass[UTF8, fontset=fandol, a4paper, 12pt]{ctexart}
\usepackage[margin=2.5cm]{geometry}
\usepackage{amsmath,amssymb,amsthm}
\usepackage{enumitem,fancyhdr}
\usepackage[hidelinks]{hyperref}
\newcommand{\R}{\mathbb R}
\newcommand{\Nplus}{\mathbb N_{+}}
\theoremstyle{definition}
\newtheorem{definition}{定义}
\newtheorem{theorem}{定理}
\renewcommand{\thetheorem}{2.\arabic{theorem}}
\renewcommand{\proofname}{证明}
\setlist[enumerate]{leftmargin=2em,itemsep=0.2em,topsep=0.3em}
\setlength{\headheight}{16pt}
\pagestyle{fancy}
\fancyhf{}
\fancyhead[L]{\small 数学分析（一）}
\fancyhead[R]{\small 2.2\quad 收敛数列的性质}
\fancyfoot[C]{\thepage}
\renewcommand{\headrulewidth}{0.3pt}
\raggedbottom
\title{收敛数列的性质}
\author{Yuchen Zhang}
\date{整理日期：2026年9月23日}
\makeatletter
\renewcommand{\maketitle}{\begin{center}
{\LARGE\bfseries \@title\par}\vspace{0.6em}
\@author\par\vspace{0.2em}{\small\@date\par}
\end{center}}
\makeatother
\begin{document}
\maketitle
\setcounter{section}{1}
\setcounter{theorem}{1}
\section{收敛数列的性质}

以下均设 $\{x_n\}$ 为实数列。若存在 $a\in\R$ 使
\[
  \forall\varepsilon>0,\quad \exists N\in\Nplus,\quad
  \forall n>N,\quad |x_n-a|<\varepsilon,
\]
则记 $x_n\to a$。

\begin{theorem}[唯一性]
若 $x_n\to a$ 且 $x_n\to b$，则 $a=b$。
\end{theorem}
\begin{proof}
反设 $a\ne b$。取 $\varepsilon=|a-b|/3>0$。由 $x_n\to a$，存在
$N_1$ 使 $n>N_1$ 时 $|x_n-a|<\varepsilon$；由 $x_n\to b$，存在
$N_2$ 使 $n>N_2$ 时 $|x_n-b|<\varepsilon$。
取 $n>\max\{N_1,N_2\}$，由三角不等式得
\[
 |a-b|\le |a-x_n|+|x_n-b|<2\varepsilon
 =\frac23|a-b|,
\]
矛盾。因此 $a=b$。
\end{proof}

\begin{theorem}[有界性]
若 $x_n\to a$，则数列 $\{x_n\}$ 有界，即存在 $M>0$，使得对一切
$n\in\Nplus$ 都有 $|x_n|\le M$。
\end{theorem}
\begin{proof}
取 $\varepsilon=1$。由 $x_n\to a$，存在 $N\in\Nplus$，使 $n>N$ 时
\[
 |x_n-a|<1.
\]
于是对 $n>N$，
\[
 |x_n|\le |x_n-a|+|a|<1+|a|.
\]
前 $N$ 项只有有限项，令
\[
 M=\max\bigl\{|x_1|,\ldots,|x_N|,1+|a|\bigr\},
\]
则对所有 $n\in\Nplus$ 都有 $|x_n|\le M$，故数列有界。
\end{proof}

\newpage
\begin{theorem}[保号性]
若 $x_n\to a$ 且 $a>0$，则存在 $N\in\Nplus$，使 $n>N$ 时 $x_n>0$。
若 $a<0$，则存在 $N\in\Nplus$，使 $n>N$ 时 $x_n<0$。
\end{theorem}
\begin{proof}
先设 $a>0$。在极限定义中取 $\varepsilon=a/2$，则存在 $N$，使 $n>N$ 时
\[
 |x_n-a|<\frac a2.
\]
因此
\[
 x_n\ge a-|x_n-a|>a-\frac a2=\frac a2>0.
\]
若 $a<0$，取 $\varepsilon=|a|/2$，则
\[
 x_n\le a+|x_n-a|<a+\frac{|a|}{2}=\frac a2<0
\]
对充分大的 $n$ 成立。
\end{proof}

\subsection*{最终区间性质}
保号性可推广为：若 $x_n\to a$，且 $b<a<c$，则存在 $N\in\Nplus$，使
\[
 n>N\quad\Longrightarrow\quad b<x_n<c.
\]
证明时取
\[
 \varepsilon=\min\{a-b,c-a\}>0.
\]
当 $n>N$ 时，$|x_n-a|<\varepsilon$，于是
\[
 b\le a-\varepsilon<x_n<a+\varepsilon\le c.
\]
因此，收敛数列从某一项起始终落在任意包含其极限的开区间内。
\end{document}
